\section*{Erd\H{o}s Problem 1163 --- GPT 6 Astra Ultra partial attempt}


\subsection*{FORMAL RESTATEMENT}
Study, using statistical descriptions, the arithmetic properties of the orders of subgroups $|H|$ with $H\leq S_n$ as $n$ varies.

\subsection*{CONTEXT AND SCOPE}
This replaces the earlier statement-only record with an elementary mathematical attempt. The arguments below establish only their stated partial results; no novelty or full solution is claimed.

The official page explicitly calls the intended meaning unclear. It does not choose between statistics of distinct attainable orders, a uniformly sampled subgroup, or a randomly generated subgroup. This attempt preserves that ambiguity.

\subsection*{ATTACK PLAN AND WORK}
To expose the statistical-model issue quantitatively, consider the concrete
family $E\cong\mathbb F_2^k\le S_{2k}$ generated by disjoint transpositions.
It has exactly $k+1$ distinct subgroup orders, namely $2^j$ for $0\le j\le k$.
Uniform sampling of these orders makes $j$ uniform. Uniform sampling of its
subgroups gives a quite different distribution.

Counting ordered bases of subspaces gives
\[N_{k,j}=\prod_{i=0}^{j-1}\frac{2^k-2^i}{2^j-2^i},\qquad
2^{j(k-j)}\le N_{k,j}\le4\,2^{j(k-j)}.\tag{1}\]
For the upper bound, the denominator after normalization is
$\prod_{h=1}^j(1-2^{-h})\ge1/4$: the first two factors give $3/8$, and the
remaining product is at least $1-\sum_{h\ge3}2^{-h}=3/4$.
The lower bound follows by comparing each normalized numerator factor with
the corresponding denominator factor.

Since there are at least $2^{\lfloor k^2/4\rfloor}$ subgroups in total,
uniform subgroup sampling satisfies
\[\Pr(|j-k/2|\ge t)\le4(k+1)2^{-t^2+1/4}.\tag{2}\]
Thus this distribution is concentrated near the middle rank, unlike uniform
sampling of the distinct attainable orders. Equations (1)--(2) concern subgroups
of $E$, not all subgroups of $S_{2k}$; that restriction is essential.

\subsection*{VERIFICATION AND REMAINING GAP}
Fix the intended sampling rule, then analyze the full symmetric-group
subgroup family. The example proves that two natural interpretations need not
have the same answer; it does not choose the intended interpretation for the source.

\subsection*{SOURCES}
https://www.erdosproblems.com/1163

https://www.erdosproblems.com/forum/thread/1163

\subsection*{FINAL}
\textbf{UNRESOLVED.} The full source problem remains outside the scope of the proved partial results.

\subsection*{COMPLETION ESTIMATE}
$5\%$. This is a conservative subjective estimate toward the whole problem, not a probability that the argument is correct or a claim that a fixed fraction of the problem has been solved.
